\documentclass[10pt,a4paper]{amsart}
\usepackage[margin=19mm]{geometry}
\usepackage{amsmath,amssymb,mathtools}
\usepackage[scr=rsfs,cal=euler]{mathalfa}
\usepackage{xfrac}
\usepackage[unicode,hidelinks,bookmarks=false]{hyperref}
\newcommand{\ie}{i.e.\ }
\newcommand{\cf}{cf.\ }
\newcommand{\resp}{respectively\ }
\newcommand{\Subsection}{\subsection}
\providecommand{\nobreakdash}{\nobreak\hbox{-}}
\setlength{\emergencystretch}{2em}
\multlinegap0pt
\usepackage{amssymb}
\usepackage{xcolor}
\usepackage{enumerate}
\usepackage{bm}
\usepackage{dsfont}
\usepackage{mathtools}
\newtheorem{proposition}{Proposition}
\title{Explicit constructions of anti-automorphisms of cyclic and generalized cyclic algebras}
\author{Susanne Pumpluen}
\date{Source: arXiv:2603.27291v2 (CC BY 4.0)}
\begin{document}
\maketitle

\noindent\emph{Source and licence.} Explicit constructions of anti-automorphisms of cyclic and generalized cyclic algebras. Version arXiv:2603.27291v2 (CC BY 4.0), distributed by arXiv under the Creative Commons Attribution 4.0 International licence, http://creativecommons.org/licenses/by/4.0/, as recorded in the arXiv licence metadata for this exact version and shown on the abstract page https://arxiv.org/abs/2603.27291v2. Full licence notice: \texttt{licenses/arxiv-2603.27291-CC-BY-4.0.txt}.\par\smallskip

\noindent\textit{Editorial scope.} Counted unit: the article's proposition prop:newconditiononk (source0.tex lines 657--661). The article's complete original proof of it is delivered with it (source0.tex lines 663--692, one \texttt{proof} environment of 30 lines). No mathematics is written, repaired or re-derived: the statement and the proof are the article's own lines, in the article's own notation, and no retained line is substituted or re-flowed.  No further source-local context is needed: every cross-reference of every retained line resolves inside the two delivered spans.  Nothing was omitted from any retained span, so the package carries no omission record.  No retained line contains a citation, so the wrapper carries no bibliography and no bibliography entry.  No retained line contains a figure environment, an image-inclusion command or drawing code, so no image is delivered and the package declares no asset dependency.  Editorial formatting only, in addition: an AMS \texttt{amsart} preamble that restates the article's own declarations and macros used by the retained lines, namely the environments \texttt{proposition} on separate counters, together with the standard packages those lines need.  The retained environments print in this document as Proposition 1 in order of appearance, whereas the article prints the same items with the numbering of its own full document.  The complete native source of the article as distributed in the arXiv e-print for this exact version is bundled unchanged as \texttt{sources/197336/source0.tex}.
 \begin{proposition}\label{prop:newconditiononk}
   Suppose that $\tau \sigma\tau^{-1}=\sigma^k$ for some $k\in \{2,\dots, n-1 \}$ with $\gcd (k, n)=1$ and suppose that $n> m$.
   Then there is no ring isomorphism
    $G:K[t;\sigma]/K[t;\sigma](t^m-a)\to K[t;\sigma]/K[t;\sigma](t^m-b)$ extending $\tau$.
 \end{proposition}

 \begin{proof}
 Suppose  that $G(t) = \sum_{i=0}^{m-1} \alpha_i t^i$ for some $\alpha_i \in K$.
Then we have
\begin{equation} \label{eqn:automorphism_necessity_theorem 1}
G(tz) = G(t)G(z) = (\sum_{i=0}^{m-1} \alpha_i t^i ) \tau(z) =
\sum_{i=0}^{m-1} \alpha_i \sigma^{i}(\tau(z)) t^i \end{equation} and
\begin{equation} \label{eqn:automorphism_necessity_theorem 2}
G(tz) =
G(\sigma(z)t) = \sum_{i=0}^{m-1} \tau(\sigma(z)) \alpha_i t^i
\end{equation}
for all $z \in K$. Comparing the coefficients of $t^i$
in \eqref{eqn:automorphism_necessity_theorem 1} and
\eqref{eqn:automorphism_necessity_theorem 2} we obtain
$$\alpha_i \sigma^{i}(\tau(z)) = \tau(\sigma(z)) \alpha_i$$
for all $i \in\{ 0, \ldots, n-1\}$ and all $z \in K$. This is equivalent to
$$\alpha_i (\sigma^i(\tau(z)) - \tau\sigma(z)  )=0,$$
for all $i \in\{ 0, \ldots, m-1 \}$ and all $z \in K$.
Since $k\in \{2,\dots, n-1 \}$, $\gcd (k, n)=1$,
 $\tau \sigma\tau^{-1}=\sigma^k$
 and $\sigma$ has order $n > m$, we know that
 $$\tau \sigma\tau^{-1}\not=\sigma^i$$
 for all  $i \in\{ 0, \ldots, m-1\}$ with $i\not=k$ so that  $\alpha_i = 0$ for
all $ i \in\{ 0, \ldots, m-1\}$  with $i\not=k$.

 If  $k\geq m$ then this implies that $G(t)=0$ and so $G$ cannot be an isomorphism. This implies we need $k<m$.
If $k<m$ then this implies that $G(t) = \alpha t^k$ for some $\alpha\in K^{\times}$.
We obtain $G=G_{\tau,\alpha,k}$
since $G|_{K}=\tau$ and $G(t)=\alpha t^k$.
However, the map $G=G_{\tau,\alpha,k}$ is only an isomorphism when $n\mid m$, so this cannot happen.
\end{proof}

\end{document}
